\documentclass[a4paper,11pt]{ltjsarticle}
\usepackage[haranoaji]{luatexja-preset}
\usepackage{amsmath,amssymb}
\usepackage[top=12mm,bottom=10mm,left=14mm,right=14mm]{geometry}
\usepackage{array}
\usepackage{tikz}
\usetikzlibrary{calc}
\pagestyle{empty}
\setlength{\parindent}{0pt}
\newlength{\qcol}\setlength{\qcol}{0.47\textwidth}
\newlength{\qgap}
\newlength{\anscolw}\setlength{\anscolw}{0.30\textwidth}
\newlength{\ansh}
\newcommand{\Q}[2]{\parbox[t]{\qcol}{\raggedright (#1)\hspace{0.6em}$\displaystyle #2$}}
\newcommand{\QR}[4]{\Q{#1}{#2}\hfill\Q{#3}{#4}\par\vspace{\qgap}}
\newcommand{\anscell}[1]{\rule[-\ansdrop]{0pt}{\ansh}(#1)}
\newlength{\ansdrop}

\begin{document}
{\large\bfseries 計算問題　17}\hfill 名前(\underline{\hspace{32mm}})\par\vspace{3mm}
■（1）〜（16）の計算をしなさい。（17）〜（20）は方程式を解きなさい。\par\vspace{3mm}
\setlength{\qgap}{5.4mm}
\QR{1}{(-9)+(-16)-(-25)}{2}{(-8)-(-13)}
\QR{3}{\dfrac{x+6}{3}\times9}{4}{(36a-6)\div6}
\QR{5}{\dfrac{5}{4}(8x+4)+\dfrac{3}{2}(4x+2)}{6}{-2^4+(-3)^2+(-5)}
\QR{7}{-3a\times(-5)}{8}{-24x\div6}
\QR{9}{(-11)\times(-111)}{10}{\dfrac{1}{3}(9a-12)-\dfrac{1}{2}(8a-6)}
\QR{11}{(5a+1)+(2a-3)}{12}{(2x+7)-(5x-6)}
\QR{13}{\dfrac{4}{3}\div\dfrac{2}{5}\times(-24)}{14}{7-(-9)\div\left(-\dfrac{3}{4}\right)}
\QR{15}{8-12\div(-3)-(-5)}{16}{2x+3-\{4x-(5x+6)\}}
\QR{17}{-3x+5=-7}{18}{5x-4=3x-6}
\QR{19}{0.5x+0.6=0.2x}{20}{\dfrac{2}{5}(x+2)=\dfrac{1}{4}x+2}
\vspace{1mm}
\Q{21}{\left\{\begin{array}{l} y=2x+3 \\ y=-2x+5 \end{array}\right.}\par\vspace{3mm}
\setlength{\ansh}{9.5mm}\setlength{\ansdrop}{6mm}%
\begin{center}\begin{tabular}{|p{\anscolw}|p{\anscolw}|p{\anscolw}|}\hline
\anscell{1} & \anscell{2} & \anscell{3} \\\hline
\anscell{4} & \anscell{5} & \anscell{6} \\\hline
\anscell{7} & \anscell{8} & \anscell{9} \\\hline
\anscell{10} & \anscell{11} & \anscell{12} \\\hline
\anscell{13} & \anscell{14} & \anscell{15} \\\hline
\anscell{16} & \anscell{17} & \anscell{18} \\\hline
\anscell{19} & \anscell{20} & \anscell{21} \\\hline
\end{tabular}\end{center}

\end{document}
